\documentclass{article}
\usepackage{csquotes}
\usepackage[a4paper]{geometry}
\usepackage{fancyhdr}
\pagestyle{fancy}
\lhead{Die EU}
\rhead{Februar 2026}
\begin{document}
\section{Die EU}
Auch die EU sorgt sich teilweise um ihre Sicherheit. Dies tut sie unter anderem durch die GASP.
 
\subsection{Die GASP} 
Schon seit 1992 gibt es die \emph{GASP}, die Gemeinsame Außen- und Sicherheitspolitik. Der Europäische Rat muss dessen Leitlinien einstimmig beschließen bevor die AußenministerInnen im Rat der EU die verbindlichen Maßnahmen bestimmen.
 
Die können in drei Arten aufgeteilt werden, 
\begin{enumerate}
 \item Aktionen der EU, früher, und passender, \textquote{Gemeinsame Aktionen}. Hierzu zählen z.\,B. Sanktionen.
 \item Gemeinsame Standpunkte, also Leitlinien, für die Außenpolitik der einzelnen Mitgliedsstaaten.
 \item Duchführungsbeschlüsse, welche die Umsetzmodalitäten der Standpunkte regeln.
\end{enumerate} 
Aufgrund der geforderten Einstimmigkeit ist der gefundene Konsens minimal. 
 
Beispielsweise konnte bereits eine Einstimmigkeit für Sanktionen gegen Russland, keiner aber für Waffenlieferungen die Ukraine gefunden werden. 
\end{document}